% GENERATED FILE. Edit canonical lessons, then run npm run generate:books.
% canonical-source-hash: fnv1a64:82aa6d0040ccd9eb
% canonical-lessons: JA-C28-tsukau, JA-C28-utau, JA-C28-iu, JA-C28-omou, JA-C28-motsu

\chapter{To use, To sing, To say, To think, To hold}
\label{ch:ja-to-use-to-sing-to-say-to-think-to-hold}
\hlchaptermodality{28}

\begin{chapteropening}
\textbf{By the end of this chapter:} \emph{I can say to use, to sing, to say, to think and to hold.}

Complete the last lesson of chapter 28: I can say to use, to sing, to say, to think and to hold.
\end{chapteropening}

\section[\texorpdfstring{\ja{つかう} (tsukau)}{つかう (tsukau)}]{\ja{つかう} (tsukau) — to use}

\label{lesson:JA-C28-tsukau}



\begin{quote}
Before the new one: say the Japanese for to wait, then the Japanese for to meet.
\end{quote}

\subsection*{You'll want to know: \ja{つかう}}
\textbf{\ja{つかう}} — \emph{tsukau} — \textquotedblleft{}to use\textquotedblright{}.

\begin{grammarlens}[title={What you've built}]
One more word for this chapter.
\end{grammarlens}

\subsection*{Your turn}
\begin{itemize}
\raggedright
  \item \emph{Say it:} \emph{tsukau}
  \item \emph{Say it:} \emph{tsukau}, once more
  \item \emph{Say it:} say \emph{tsukau}
  \item \emph{Recall:} say the Japanese for to wait, then the Japanese for to meet, then say \emph{tsukau} again
\end{itemize}

\begin{tcolorbox}[breakable,colback=teal!4,colframe=teal!35!black,title={Before you move on}]
What does \ja{つかう} mean? (\textquotedblleft{}To use\textquotedblright{}.) Say it once more.
\end{tcolorbox}
\section[\texorpdfstring{\ja{うたう} (utau)}{うたう (utau)}]{\ja{うたう} (utau) — to sing}

\label{lesson:JA-C28-utau}



\begin{quote}
Before the new one: say the Japanese for to meet, then the Japanese for to use.
\end{quote}

\subsection*{You'll want to know: \ja{うたう}}
\textbf{\ja{うたう}} — \emph{utau} — \textquotedblleft{}to sing\textquotedblright{}.

\begin{grammarlens}[title={What you've built}]
One more word for this chapter.
\end{grammarlens}

\subsection*{Your turn}
\begin{itemize}
\raggedright
  \item \emph{Say it:} \emph{utau}
  \item \emph{Say it:} \emph{utau}, once more
  \item \emph{Say it:} say \emph{utau}
  \item \emph{Recall:} say the Japanese for to meet, then the Japanese for to use, then say \emph{utau} again
\end{itemize}

\begin{tcolorbox}[breakable,colback=teal!4,colframe=teal!35!black,title={Before you move on}]
What does \ja{うたう} mean? (\textquotedblleft{}To sing\textquotedblright{}.) Say it once more.
\end{tcolorbox}
\section[\texorpdfstring{\ja{言}\ja{う} (iu)}{言う (iu)}]{\ja{言}\ja{う} (iu) — to say}

\label{lesson:JA-C28-iu}



\begin{quote}
Before the new one: say the Japanese for to use, then the Japanese for to sing.
\end{quote}

\subsection*{You'll want to know: \ja{言}\ja{う}}
\textbf{\ja{言}\ja{う}} (\ja{いう}) — \emph{iu} — \textquotedblleft{}to say\textquotedblright{}.

Its first sign is \textbf{\ja{言}}, the speech side of \textbf{\ja{語}} that you have written since chapter five; the \textbf{\ja{う}} after it is hiragana.

\begin{grammarlens}[title={What you've built}]
One more word for this chapter.
\end{grammarlens}

\subsection*{Your turn}
\begin{itemize}
\raggedright
  \item \emph{Say it:} \emph{iu}
  \item \emph{Read it:} \textbf{\ja{言}\ja{う}}, then say \emph{iu} once more
  \item \emph{Say it:} say \emph{iu}
  \item \emph{Recall:} say the Japanese for to use, then the Japanese for to sing, then say \emph{iu} again
\end{itemize}

\begin{tcolorbox}[breakable,colback=teal!4,colframe=teal!35!black,title={Before you move on}]
What does \ja{言}\ja{う} mean? (\textquotedblleft{}To say\textquotedblright{}.) Say it once more.
\end{tcolorbox}
\section[\texorpdfstring{\ja{おもう} (omou)}{おもう (omou)}]{\ja{おもう} (omou) — to think}

\label{lesson:JA-C28-omou}



\begin{quote}
Before the new one: say the Japanese for to sing, then the Japanese for to say.
\end{quote}

\subsection*{You'll want to know: \ja{おもう}}
\textbf{\ja{おもう}} — \emph{omou} — \textquotedblleft{}to think\textquotedblright{}.

\begin{grammarlens}[title={What you've built}]
One more word for this chapter.
\end{grammarlens}

\subsection*{Your turn}
\begin{itemize}
\raggedright
  \item \emph{Say it:} \emph{omou}
  \item \emph{Say it:} \emph{omou}, once more
  \item \emph{Say it:} say \emph{omou}
  \item \emph{Recall:} say the Japanese for to sing, then the Japanese for to say, then say \emph{omou} again
\end{itemize}

\begin{tcolorbox}[breakable,colback=teal!4,colframe=teal!35!black,title={Before you move on}]
What does \ja{おもう} mean? (\textquotedblleft{}To think\textquotedblright{}.) Say it once more.
\end{tcolorbox}
\section[\texorpdfstring{\ja{もつ} (motsu)}{もつ (motsu)}]{\ja{もつ} (motsu) — to hold}

\label{lesson:JA-C28-motsu}



\begin{quote}
Before the new one: say the Japanese for to say, then the Japanese for to think.
\end{quote}

\subsection*{You'll want to know: \ja{もつ}}
\textbf{\ja{もつ}} — \emph{motsu} — \textquotedblleft{}to hold\textquotedblright{}.

\begin{grammarlens}[title={What you've built}]
One more word for this chapter.
\end{grammarlens}

\subsection*{Your turn}
\begin{itemize}
\raggedright
  \item \emph{Say it:} \emph{motsu}
  \item \emph{Say it:} \emph{motsu}, once more
  \item \emph{Say it:} say \emph{motsu}
  \item \emph{Recall:} say the Japanese for to say, then the Japanese for to think, then say \emph{motsu} again
\end{itemize}

\begin{tcolorbox}[breakable,colback=teal!4,colframe=teal!35!black,title={Before you move on}]
What does \ja{もつ} mean? (\textquotedblleft{}To hold\textquotedblright{}.) Say it once more.
\end{tcolorbox}
